% !TEX root = ../../hw01.tex
% Homework 01, Exercise 6. Shared solution.

% Completed from the argument on the right page of IMG_8999.
\begin{proof}[Solution]
Let $\mathbb{A}$ be the set of real algebraic numbers and
$\mathbb{T}$ the set of real transcendental numbers. An algebraic
number is a root of a nonzero polynomial with integer coefficients;
a transcendental number is not such a root. Hence
\[
  \R=\mathbb{A}\cup\mathbb{T},
  \qquad\mathbb{A}\cap\mathbb{T}=\varnothing.
\]
First, we show that $\mathbb{A}$ is countable. For each $n\ge1$,
consider the nonzero polynomials
\[
  P(X)=c_0+c_1X+\cdots+c_nX^n,
  \qquad c_j\in\Z,\quad |c_j|\le n.
\]
Each of the $n+1$ coefficients has $2n+1$ possible values, so
there are at most $(2n+1)^{n+1}$ such polynomials.
A nonzero polynomial of degree at most $n$ has at most $n$
distinct real roots. Thus the set $F_n$ of all their real roots
is finite.

Every algebraic number is a root of a nonzero integer polynomial.
Taking $n$ at least as large as its degree and the absolute values
of all its coefficients places that root in $F_n$.
Conversely, every root in $F_n$ is algebraic. Hence
\[
  \mathbb{A}=\bigcup_{n=1}^{\infty}F_n.
\]
Within each finite set $F_n$, list the roots in increasing order.
Read these finite lists in the order $F_1,F_2,\ldots$, skipping
repeated roots. Each root is reached after finitely many steps.
This proves that $\mathbb{A}$ is countable. Every $n\in\N$
is a root of $X-n$, so $\N\subseteq\mathbb{A}$. Thus
$\mathbb{A}$ is infinite, and $|\mathbb{A}|=\aleph_0$.

\begin{marginfigure}
  \centering
  % !TEX root = ../../problems/hw01.tex
\begin{tikzpicture}[x=1cm,y=1cm,>=Stealth,font=\footnotesize]
  \draw[->,black!55] (0,0) -- (3.8,0);
  \draw[black!65,fill=white] (0.25,0) circle[radius=2.5pt];
  \fill (3.45,0) circle[radius=3pt];
  \node[below=6pt] at (0.25,0) {$\aleph_0$};
  \node[below=6pt] at (3.15,0) {$|\mathbb{T}|=\aleph$};
  \node[above=8pt,align=center] at (1.85,0)
    {No intermediate\\cardinality under CH};
\end{tikzpicture}

  \caption[The continuum hypothesis.]{CH excludes intermediate cardinalities. Since
    $\mathbb{T}$ is uncountable and contained in $\R$, its cardinality
    must equal $\aleph$.}
  \label{fig:hw01:continuum-gap}
\end{marginfigure}

Now suppose that $\mathbb{T}$ were countable. Then
$\R=\mathbb{A}\cup\mathbb{T}$ would be a union of two countable
sets and therefore countable. This contradicts the uncountability
of $\R$. Hence $\mathbb{T}$ is uncountable. Since
$\mathbb{T}\subseteq\R$, we have
\[
  \aleph_0<|\mathbb{T}|\le|\R|=\aleph.
\]
By CH there is no cardinality strictly between $\aleph_0$ and
$\aleph$, as illustrated in Figure~\ref{fig:hw01:continuum-gap}.
The lower endpoint is already excluded, so the only possibility is
\[
  |\mathbb{T}|=\aleph.\qedhere
\]
\end{proof}
